\chapter{Hermitian Eigensystems and Spectral Subspaces}
\index{eigensystem|textbf}
\index{Hermitian matrix|textbf}
\index{spectral theorem|textbf}
\index{eigenspace|textbf}

\label{ch:training-hermitian-eigensystems}

\chapterstatus{pedagogical finite-dimensional foundation. All matrices and
numerical examples are synthetic. This chapter establishes no material result,
discretization-convergence claim, scientific validation, or uncertainty
quantification.}

\section{Learning goals}

After working through this chapter, a reader should be able to

\begin{enumerate}
  \item distinguish an eigenpair, spectrum, eigenspace, assembled eigensystem,
        and spectral decomposition;
  \item explain why Hermitian matrices have real eigenvalues and orthogonal
        eigenspaces;
  \item derive a small characteristic equation without mistaking it for a
        production numerical algorithm;
  \item interpret degeneracy through invariant subspaces and spectral
        projectors rather than arbitrary eigenvector columns;
  \item use the Rayleigh quotient to characterize extreme eigenvalues; and
  \item verify a computed Hermitian eigensystem with residual, orthogonality,
        and reconstruction diagnostics.
\end{enumerate}

Chapter~\ref{ch:training-finite-representations} distinguished a linear map
from its matrix in a selected basis. We now ask when a basis can be chosen so
that a matrix acts by separately scaling its basis vectors. That question is
central to the monograph: finite Hamiltonians are diagonalized, low-energy
spaces are retained, Bloch bands are assembled from fiber eigenvalues, and
composite frames are compared through their subspaces rather than through
arbitrary raw columns.

\section{Eigenpairs are invariant-line statements}
\index{eigenvalue}
\index{eigenvector}
\index{spectrum}

Let
\begin{equation}
  H\in\mathbb C^{N\times N}.
\end{equation}
An \emph{eigenpair} consists of a scalar \(\lambda\in\mathbb C\) and a nonzero
vector \(\mathbf v\in\mathbb C^N\) satisfying
\begin{equation}
  H\mathbf v=\lambda\mathbf v,
  \qquad
  \mathbf v\neq\mathbf 0.
  \label{eq:training-eigenpair}
\end{equation}
The line spanned by \(\mathbf v\) is invariant under \(H\): applying \(H\)
keeps the result on that line. Saying that the ``direction does not change'' is
less precise because a negative real eigenvalue reverses an oriented vector,
and a complex eigenvalue changes its phase.

Several related terms should not be conflated.

\begin{description}
  \item[Eigenvalue.] One scalar \(\lambda\) satisfying
        Equation~\eqref{eq:training-eigenpair} for at least one nonzero vector.
  \item[Eigenvector.] One nonzero vector associated with that eigenvalue.
  \item[Eigenspace.] The complete null space
        \(\ker(H-\lambda I_N)\), including the zero vector.
  \item[Spectrum.] The collection of eigenvalues, counted with a declared
        multiplicity convention when needed.
  \item[Assembled eigensystem.] Eigenvalues and selected eigenvector columns,
        together with their ordering and normalization conventions.
  \item[Spectral decomposition.] A representation of the matrix in terms of
        its eigenvalues and invariant spectral projectors.
\end{description}

A general square matrix need not have an orthonormal eigenbasis and may even be
defective. The remainder of this chapter concerns Hermitian matrices, for
which a much stronger theorem is available.

\section{Why Hermiticity changes the problem}
\index{Hermitian matrix}
\index{self-adjoint matrix}

A complex matrix is Hermitian when
\begin{equation}
  H=H^{\dagger},
  \label{eq:training-hermitian-definition}
\end{equation}
where \(\dagger\) denotes conjugate transpose. A real symmetric matrix,
\(H=H^{\mathsf T}\), is the real-valued special case.

Suppose \((\lambda,\mathbf v)\) is an eigenpair of a Hermitian matrix. Using an
inner product that is conjugate-linear in its first argument,
\begin{equation}
  \langle\mathbf v,H\mathbf v\rangle
  =\lambda\langle\mathbf v,\mathbf v\rangle.
\end{equation}
Hermiticity also implies that
\(\langle\mathbf v,H\mathbf v\rangle\) is real. Because
\(\langle\mathbf v,\mathbf v\rangle>0\), the eigenvalue \(\lambda\) must be
real.

Now let \(\mathbf v_m\) and \(\mathbf v_n\) have distinct eigenvalues. Then
\begin{align}
  \lambda_n\langle\mathbf v_m,\mathbf v_n\rangle
  &=\langle\mathbf v_m,H\mathbf v_n\rangle,\\
  \lambda_m\langle\mathbf v_m,\mathbf v_n\rangle
  &=\langle H\mathbf v_m,\mathbf v_n\rangle.
\end{align}
The two right-hand sides are equal by Hermiticity, so
\begin{equation}
  (\lambda_n-\lambda_m)
  \langle\mathbf v_m,\mathbf v_n\rangle=0.
\end{equation}
Distinct eigenvalues therefore have orthogonal eigenspaces.

The finite-dimensional spectral theorem adds the crucial completeness result:
every Hermitian \(N\times N\) matrix has \(N\) orthonormal eigenvectors when
multiplicity is counted \parencite{hornJohnson2013}. The reality and
orthogonality arguments above explain
important parts of the theorem, but they do not by themselves prove the
existence of a complete eigenbasis.

\section{From individual eigenpairs to a spectral decomposition}
\index{eigendecomposition}

Choose normalized eigenvectors and place them in columns,
\begin{equation}
  V=
  \begin{bmatrix}
    \mathbf v_1 & \mathbf v_2 & \cdots & \mathbf v_N
  \end{bmatrix},
  \qquad
  V^{\dagger}V=VV^{\dagger}=I_N.
  \label{eq:training-eigenvector-matrix}
\end{equation}
Place the correspondingly ordered eigenvalues on a diagonal matrix,
\begin{equation}
  \Lambda
  =\operatorname{diag}(\lambda_1,\ldots,\lambda_N).
  \label{eq:training-eigenvalue-matrix}
\end{equation}
The individual equations \(H\mathbf v_n=\lambda_n\mathbf v_n\) then become
\begin{equation}
  HV=V\Lambda.
  \label{eq:training-assembled-eigensystem}
\end{equation}
Multiplying by \(V^{\dagger}\) gives the equivalent forms
\begin{equation}
  H=V\Lambda V^{\dagger},
  \qquad
  V^{\dagger}HV=\Lambda.
  \label{eq:training-hermitian-diagonalization}
\end{equation}
The first reconstructs the matrix in the original basis. The second represents
the same linear map in the eigenvector basis. Diagonalization is therefore a
basis transformation with a basis whose existence follows from Hermiticity.

Define the rank-one projector
\begin{equation}
  P_n=\mathbf v_n\mathbf v_n^{\dagger}.
\end{equation}
Then
\begin{equation}
  I_N=\sum_{n=1}^{N}P_n,
  \qquad
  H=\sum_{n=1}^{N}\lambda_nP_n.
  \label{eq:training-rank-one-spectral-decomposition}
\end{equation}
For any vector
\begin{equation}
  \mathbf x=\sum_{n=1}^{N}c_n\mathbf v_n,
  \qquad
  c_n=\mathbf v_n^{\dagger}\mathbf x,
\end{equation}
the action of \(H\) is
\begin{equation}
  H\mathbf x=\sum_{n=1}^{N}\lambda_nc_n\mathbf v_n.
\end{equation}
Thus the eigenbasis resolves a coupled matrix action into independent scalar
actions. In a degenerate eigenspace, however, the individual rank-one
projectors depend on the selected basis; their sum over the complete degenerate
space does not.

\section{The characteristic equation}
\index{characteristic equation}
\index{determinant!geometric meaning}

Move the right-hand side of Equation~\eqref{eq:training-eigenpair} to the left:
\begin{equation}
  (H-\lambda I_N)\mathbf v=\mathbf 0.
\end{equation}
A nonzero solution exists only when \(H-\lambda I_N\) is singular. Therefore
\begin{equation}
  \det(H-\lambda I_N)=0.
  \label{eq:training-matrix-characteristic-equation}
\end{equation}
This is the characteristic equation. Its left-hand side is a degree-\(N\)
polynomial in \(\lambda\), and its roots are the eigenvalues.

For a real \(N\times N\) matrix, \(|\det A|\) is its
\(N\)-dimensional volume-scaling factor. For a complex matrix viewed as a map
of a \(2N\)-dimensional real space, the corresponding real-volume factor is
\(|\det A|^2\). In either scalar field, \(\det A=0\) means that the map
collapses at least one nonzero direction and is not invertible. Applied to
\(A=H-\lambda I_N\), that collapsed direction is an eigenvector. The
determinant therefore detects the eigenvalue; the null space identifies its
eigenspace.

The characteristic polynomial is basis invariant. If
\(H'=U^{\dagger}HU\) for a unitary \(U\), then
\begin{align}
  \det(H'-\lambda I_N)
  &=\det\!\left(U^{\dagger}(H-\lambda I_N)U\right)\\
  &=\det(U^{\dagger})\det(H-\lambda I_N)\det(U)\\
  &=\det(H-\lambda I_N).
\end{align}
The entries change under the basis transformation, while the roots remain the
same.

\subsection{A complete two-dimensional example}

Consider
\begin{equation}
  H=
  \begin{bmatrix}
    2 & 1\\
    1 & 2
  \end{bmatrix}.
\end{equation}
The characteristic equation is
\begin{align}
  0
  &=\det
    \begin{bmatrix}
      2-\lambda & 1\\
      1 & 2-\lambda
    \end{bmatrix}\\
  &=(2-\lambda)^2-1\\
  &=(\lambda-1)(\lambda-3).
\end{align}
The eigenvalues are \(\lambda_1=1\) and \(\lambda_2=3\). Solving the two null
space problems and normalizing gives
\begin{equation}
  \mathbf v_1=\frac{1}{\sqrt2}
  \begin{bmatrix}1\\-1\end{bmatrix},
  \qquad
  \mathbf v_2=\frac{1}{\sqrt2}
  \begin{bmatrix}1\\1\end{bmatrix}.
\end{equation}
Consequently,
\begin{equation}
  V=\frac{1}{\sqrt2}
  \begin{bmatrix}
    1 & 1\\
    -1 & 1
  \end{bmatrix},
  \qquad
  \Lambda=
  \begin{bmatrix}
    1 & 0\\
    0 & 3
  \end{bmatrix},
\end{equation}
and direct multiplication verifies
\begin{equation}
  HV=V\Lambda,
  \qquad
  H=V\Lambda V^{\mathsf T}.
\end{equation}
The transpose appears because this example is real. The complex-Hermitian
formula uses \(V^{\dagger}\).

For any matrix, with eigenvalues counted according to algebraic multiplicity,
\begin{equation}
  \det H=\prod_{n=1}^{N}\lambda_n,
  \qquad
  \operatorname{tr}H=\sum_{n=1}^{N}\lambda_n.
\end{equation}
These identities are useful checks, but they do not replace eigenvectors or
spectral projectors.

Equation~\eqref{eq:training-matrix-characteristic-equation} is invaluable for
small hand calculations and for understanding singularity and multiplicity. It
is not the usual stable numerical route for a large matrix. Production
Hermitian eigensolvers use matrix factorizations or iterative subspace methods
rather than evaluating determinants and finding roots of an explicitly formed
characteristic polynomial \parencite{golubVanLoan2013}.

\section{The Rayleigh quotient and extreme eigenvalues}
\index{Rayleigh quotient}

For any nonzero vector \(\mathbf x\), define
\begin{equation}
  \rho_H(\mathbf x)
  =\frac{\mathbf x^{\dagger}H\mathbf x}
         {\mathbf x^{\dagger}\mathbf x}.
  \label{eq:training-rayleigh-quotient}
\end{equation}
This quotient is real when \(H\) is Hermitian. Expand
\(\mathbf x=\sum_n c_n\mathbf v_n\). The spectral decomposition gives
\begin{equation}
  \rho_H(\mathbf x)
  =\frac{\sum_n|c_n|^2\lambda_n}{\sum_n|c_n|^2}.
  \label{eq:training-rayleigh-weighted-average}
\end{equation}
The quotient is therefore a weighted average of eigenvalues. If they are
ordered as
\(\lambda_1\leq\cdots\leq\lambda_N\), then
\begin{equation}
  \lambda_1
  \leq\rho_H(\mathbf x)
  \leq\lambda_N.
\end{equation}
Moreover,
\begin{equation}
  \lambda_1=\min_{\mathbf x\neq\mathbf0}\rho_H(\mathbf x),
  \qquad
  \lambda_N=\max_{\mathbf x\neq\mathbf0}\rho_H(\mathbf x).
  \label{eq:training-rayleigh-extrema}
\end{equation}
The minimizing and maximizing vectors lie in the corresponding eigenspaces.
Intermediate eigenvalues admit related constrained variational
characterizations \parencite{hornJohnson2013}. This is the finite-dimensional
foundation of methods that seek low-energy states without computing a complete
spectrum.

\section{Degeneracy, eigenspaces, and projectors}
\index{eigenvalue!degenerate}
\index{spectral projector}

For an eigenvalue \(\lambda\), the eigenspace is
\begin{equation}
  \mathcal E_{\lambda}
  =\ker(H-\lambda I_N).
  \label{eq:training-eigenspace}
\end{equation}
For a Hermitian matrix, the dimension of this space equals the algebraic
multiplicity of \(\lambda\). A repeated eigenvalue is called degenerate when
this multiplicity exceeds one.

Let
\begin{equation}
  W_{\lambda}
  =\begin{bmatrix}
    \mathbf w_1 & \cdots & \mathbf w_g
  \end{bmatrix}
\end{equation}
contain any orthonormal basis of a \(g\)-dimensional eigenspace. Its spectral
projector is
\begin{equation}
  P_{\lambda}=W_{\lambda}W_{\lambda}^{\dagger}.
  \label{eq:training-degenerate-projector}
\end{equation}
If \(Q\in U(g)\), then \(W'_{\lambda}=W_{\lambda}Q\) is another orthonormal
basis of the same eigenspace, but
\begin{equation}
  W'_{\lambda}W_{\lambda}'^{\dagger}
  =W_{\lambda}QQ^{\dagger}W_{\lambda}^{\dagger}
  =P_{\lambda}.
\end{equation}
The basis columns are not unique; the projector is.

As a concrete example, consider
\begin{equation}
  H=
  \begin{bmatrix}
    2 & 1 & 1\\
    1 & 2 & 1\\
    1 & 1 & 2
  \end{bmatrix}.
\end{equation}
The normalized vector
\begin{equation}
  \mathbf v_+=\frac{1}{\sqrt3}
  \begin{bmatrix}1\\1\\1\end{bmatrix}
\end{equation}
has eigenvalue \(4\). Every vector orthogonal to \(\mathbf v_+\), equivalently
every \(\mathbf x\) satisfying \(x_1+x_2+x_3=0\), has eigenvalue \(1\). Thus
\begin{equation}
  \mathcal E_1
  =\left\{\mathbf x\in\mathbb R^3\mid x_1+x_2+x_3=0\right\},
  \qquad
  \dim\mathcal E_1=2,
\end{equation}
and the projector onto that plane is
\begin{equation}
  P_1=I_3-\mathbf v_+\mathbf v_+^{\mathsf T}.
\end{equation}
No particular pair of axes inside the plane is preferred by \(H\). Different
software runs or algorithms may return different valid orthonormal columns
spanning it.

Exact degeneracy must be distinguished from numerical \emph{near-degeneracy}.
Two computed eigenvalues that differ by a small amount are not automatically
mathematically equal. A calculation must declare the dimensional scale and
tolerance used to group them, and its conclusions should remain stable under
that declared rule.

\section{Phase, ordering, and subspace freedom}
\index{eigenvector!phase freedom}
\index{eigensystem!ordering}

Even a normalized nondegenerate eigenvector is not unique. For any real
\(\theta\),
\begin{equation}
  \mathbf v_n' = e^{i\theta}\mathbf v_n
\end{equation}
is an equally valid eigenvector. In a real calculation this includes the sign
change \(\mathbf v_n\mapsto-\mathbf v_n\). Eigenvalue sorting is also an output
convention rather than a theorem. Inside a degenerate eigenspace, the freedom
expands from one phase per column to an arbitrary unitary mixing of all columns
in the block.

Raw component equality is therefore not a universal comparison rule for
independently computed eigenvectors. Suitable comparisons include eigenvalue
residuals, spectral projectors, and invariant-subspace diagnostics. Later
chapters refine these ideas for momentum-dependent Bloch frames, where choices
made independently at neighboring reciprocal points affect smoothness and
real-space locality.

The spectrum is basis invariant but does not completely identify an operator.
Isospectral matrices can attach the same eigenvalues to different physical
sites, orbitals, spin channels, or spatial regions. Spectral agreement is
valuable evidence, but it supplies no physical alignment map by itself.

\section{Computing and checking a Hermitian eigensystem}
\index{eigensolver}
\index{eigenpair residual}

For a dense Hermitian matrix, use a Hermitian eigensolver such as
\texttt{numpy.linalg.eigh}, not a general nonsymmetric routine such as
\texttt{numpy.linalg.eig}. A Hermitian solver exploits the matrix structure and
returns real eigenvalues and orthonormal eigenvectors up to floating-point
error. For a large sparse matrix, computing only a requested portion of the
spectrum is often more appropriate than constructing a complete dense
factorization \parencite{golubVanLoan2013}.

A returned pair should be checked through the residual
\begin{equation}
  \mathbf r_n=H\mathbf v_n-\lambda_n\mathbf v_n.
  \label{eq:training-eigenpair-residual}
\end{equation}
For a normalized eigenvector and an energy-valued \(H\),
\(\|\mathbf r_n\|_2\) has units of energy. Its tolerance must therefore be
stated with a scale rather than treated as a dimensionless universal constant.
For several returned columns, define
\begin{equation}
  R=HV-V\Lambda.
\end{equation}
Then three distinct checks are
\begin{align}
  \|R\|_{\mathrm F}
  &\quad\text{eigenpair residual},\\
  \|V^{\dagger}V-I\|_{\mathrm F}
  &\quad\text{orthonormality defect},\\
  \|H-V\Lambda V^{\dagger}\|_{\mathrm F}
  &\quad\text{full reconstruction defect}.
\end{align}
The reconstruction check applies only when \(V\) contains a complete basis.
It should not be manufactured by densifying a large sparse operator. Partial
sparse results can be checked with sparse matrix--block products, column
orthogonality, and projectors on the returned subspace.

\begin{pythonexample}{checking eigenpairs and a degenerate subspace}
import numpy as np

# This dense matrix is synthetic and has eigenvalues 1, 1, and 4.
H = np.array(
    [[2.0, 1.0, 1.0], [1.0, 2.0, 1.0], [1.0, 1.0, 2.0]],
    dtype=np.float64,
)
np.testing.assert_allclose(H, H.T, rtol=0.0, atol=0.0)

# eigh uses Hermitian structure and returns eigenvectors as columns.
eigenvalues, eigenvectors = np.linalg.eigh(H)
Lambda = np.diag(eigenvalues)

residual = H @ eigenvectors - eigenvectors @ Lambda
orthogonality_defect = eigenvectors.T @ eigenvectors - np.eye(3)
reconstruction = eigenvectors @ Lambda @ eigenvectors.T

np.testing.assert_allclose(residual, 0.0, atol=1e-14)
np.testing.assert_allclose(orthogonality_defect, 0.0, atol=1e-14)
np.testing.assert_allclose(reconstruction, H, atol=1e-14)
np.testing.assert_allclose(eigenvalues, [1.0, 1.0, 4.0], atol=1e-14)

# Rotate the two-dimensional degenerate eigenspace without changing it.
W = eigenvectors[:, :2]
angle = 0.37
Q = np.array(
    [[np.cos(angle), -np.sin(angle)], [np.sin(angle), np.cos(angle)]],
    dtype=np.float64,
)
W_rotated = W @ Q

np.testing.assert_allclose(H @ W_rotated, W_rotated, atol=1e-14)
np.testing.assert_allclose(W_rotated @ W_rotated.T, W @ W.T, atol=1e-14)
\end{pythonexample}

The successful checks establish finite linear-algebra behavior for this
synthetic matrix. They do not establish that a discretized Hamiltonian has
converged or that any physical material is represented.

\documentcallout{Notebook to do}{Create
\path{examples/tutorials/research-monograph/foundations/hermitian_eigensystems.ipynb}
with explicit imports, explanatory inline comments, characteristic-polynomial
hand checks, Hermitian solver diagnostics, phase changes, degenerate rotations,
projector comparisons, and sparse-preserving partial-eigensystem checks. Keep
it unexecuted in the repository.}

\section{From eigenspaces to retained spaces}

This chapter treated eigenspaces selected by exact eigenvalues of one finite
Hermitian matrix. Chapter~\ref{ch:training-quantum-operators} next supplies the
quantum meaning of states, observables, and Hamiltonian eigenproblems.
Chapter~\ref{ch:training-projection-retained-spaces} later asks a different
question: which eigenvectors or spectral subspaces should be retained across a
family of represented operators, and what information is lost by that
selection? Diagonalization, projection, gauge choice, and truncation remain
separate operations.

\section{Common mistakes}

\begin{itemize}
  \item Treating \(\det(H-\lambda I)=0\) as the preferred large-scale
        numerical eigensolver.
  \item Saying an eigenvector is unique without accounting for normalization,
        phase, ordering, or degenerate rotations.
  \item Comparing raw eigenvector components from independent calculations
        without first aligning their coordinate spaces and gauge freedoms.
  \item Calling a small difference between floating-point eigenvalues an exact
        degeneracy without a declared scale and grouping tolerance.
  \item Using a full reconstruction check for a partial eigensystem.
  \item Densifying a sparse matrix merely to perform a diagnostic.
  \item Treating equality of spectra as equality of represented operators.
\end{itemize}

\section{Exercises}

\begin{enumerate}
  \item Prove directly that an eigenvalue of a Hermitian matrix is real.
  \item Derive the characteristic equation and normalized eigenvectors for
        \(\begin{bmatrix}3&i\\-i&3\end{bmatrix}\).
  \item Use Equation~\eqref{eq:training-rayleigh-weighted-average} to prove the
        bounds on the Rayleigh quotient.
  \item Construct two different orthonormal bases of the plane
        \(x_1+x_2+x_3=0\) and verify that they produce the same projector.
  \item Explain why two independently returned eigenvectors can represent the
        same nondegenerate eigendirection while failing componentwise equality.
  \item State which diagnostics remain available when only the lowest \(q<N\)
        eigenpairs of a sparse Hermitian matrix have been computed.
  \item Construct two isospectral \(3\times3\) Hermitian matrices with visibly
        different entries. State what additional alignment information is
        needed before their difference can be interpreted physically.
\end{enumerate}

\section{Evidence boundary}

Residual, orthogonality, projector, and reconstruction checks can verify a
finite eigensystem calculation within declared tolerances. They do not
establish that the parent operator is physically appropriate, that a numerical
discretization has converged, that a retained subspace is adequate for a later
observable, or that an eigensystem models a real material. Those are separate
numerical and scientific claims.
